\section{Evaluation}
\label{sec:evaluation}

The evaluation has two parts: Section~\ref{sec:eval-results} evaluates the \emph{benchmark itself}, checking that its baselines match each family's production framework so that speedups over them are meaningful; Sections~\ref{sec:eval-agents} and~\ref{sec:eval-e2e} evaluate \emph{kernel agents} on the default set, at the kernel level and end to end.

\subsection{Benchmark Validation: Baselines vs.\ Production Frameworks}
\label{sec:eval-results}

On the full set, we benchmark every architecture's end-to-end performance (throughput and latency on real workloads) and the correctness of each L1--L4 task against the state-of-the-art production framework for its family (e.g., vLLM for LLMs, SGLang for serving, diffusers for SDXL, FLA for linear-attention recurrences, timm for vision encoders).
Figure~\ref{fig:eval-e2e} reports the category-level geometric means: gains concentrate in linear attention ($1.45\times$ throughput) and recommendation ($1.51\times$), and the other categories sit near parity; Appendix~\ref{app:reference_perf} breaks the results down by architecture and GPU (Blackwell vs.\ Hopper).
References are pinned per release and GPU (vLLM~0.26.0 here); a daily CI job tracks their drift, 1--9\% for the median model over the four months before our pin (Appendix~\ref{app:drift}).


\subsection{Agent Evaluation: Kernel-Level Results}
\label{sec:eval-agents}

\paragraph{Agents.}
We evaluate Dr.~Kernel~\citep{drkernel}, an RL-trained model and one of the most recent KernelBench-ecosystem agents (others did not reproduce their results in our setting, e.g.\ KernelAgent~\citep{kernelagent}, or assume KernelBench's single-file tasks); KDA~\citep{kda} and AKO~\citep{ako} (adapted by the GEMM People entry), long-horizon harnesses behind winning MLSys 2026 FlashInfer contest entries~\citep{mlsys26contest}: KDA placed in all three full-agent tracks and GEMM People second on MoE; and Claude Code and Codex~\citep{codex}, the coding agents these harnesses build on.
Their campaigns total 6{,}900 agent-hours and 42B tokens, with Claude Code, KDA, and AKO spending a median of 6--12 hours per operator (Appendix~\ref{app:agents}).
We therefore evaluate fewer agents than single-turn leaderboards: the agents practitioners deploy iterate for hours per kernel (contest teams had 13 weeks for five kernels).
We omit evolutionary search, which multiplies this cost by population size and generation count and whose strongest systems~\citep{alphaevolve,kernelevolve} are not public.
Each agent optimizes every L1--L3 operator captured from the default set on B200 (47 L1, 64 L2, and 21 L3 targets) against the human-expert kernels that production frameworks ship.
L4 tasks are full models, so benchmarking them is the end-to-end evaluation of Section~\ref{sec:eval-e2e}.
We bench L1--L3 in three modes:
\emph{independent} (each operator generated and benched in isolation, over production lower-level kernels),
\emph{composed} (the independent winners---correct, geomean speedup $\geq 1\times$---activated together, so higher-level kernels run on the agent's own lower-level winners), and
\emph{sequential} (L2 and L3 regenerated on top of the frozen lower-level winners, benched composed).
Four agents reach geomean speedups of $1.6$--$6.6\times$ over the production kernels (Table~\ref{tab:kernel-levels}), but composition erodes them (Appendix Table~\ref{tab:kernel-composition}): once their children are also agent-generated, 26 of the agents' 71 L3 winners (37\%) fail correctness (13 of 18 for Codex), the surviving L3 winners of KDA and AKO lose 12--14\% of their speedup, and individual modules lose up to 76\% (Claude Code's \texttt{yolov10\_backbone}, $30.3\times \to 7.2\times$).
Sequential generation makes this an optimization target: it yields 80 correct L3 winners across agents instead of 45 composed, and the sequential L3 kernels of AKO and Codex are $1.48\times$ and $1.35\times$ faster than their composed independent ones.

\begin{table}[t]
  \centering
  \caption{\textbf{Kernel agents on \fastkernels{} (B200, default set).}
  \emph{Corr.}: operators correct on every benchmarked case; \emph{Geo.}: geomean speedup over correct operators; fast$_p$: correct and more than $p\times$ faster than production (Appendix~\ref{app:fastp}).
  Composed and sequential rows contain winners only.}
  \label{tab:kernel-levels}
  \small
  \setlength{\tabcolsep}{3pt}
  \resizebox{\linewidth}{!}{\begin{tabular}{@{}llrrrrrrrrrrrr@{}}
\toprule
 & & \multicolumn{4}{c}{\textbf{L1} (47 targets)} & \multicolumn{4}{c}{\textbf{L2} (64 targets)} & \multicolumn{4}{c}{\textbf{L3} (21 targets)} \\
\cmidrule(lr){3-6}\cmidrule(lr){7-10}\cmidrule(lr){11-14}
Agent & Mode & Corr. & Geo. & fast$_{1}$ & fast$_{2}$ & Corr. & Geo. & fast$_{1}$ & fast$_{2}$ & Corr. & Geo. & fast$_{1}$ & fast$_{2}$ \\
\midrule
\multirow{3}{*}{Dr.~Kernel} & Independent & 31 & 0.93$\times$ & 18 & 1 & 18 & 0.59$\times$ & 5 & 0 & 4 & 0.97$\times$ & 1 & 0 \\
 & Composed & 18 & 1.29$\times$ & 18 & 1 & 5 & 1.07$\times$ & 4 & 0 & 1 & 1.03$\times$ & 1 & 0 \\
 & Sequential & 18 & 1.29$\times$ & 18 & 1 & 5 & 1.40$\times$ & 5 & 0 & 2 & 1.16$\times$ & 2 & 0 \\
\midrule
\multirow{3}{*}{Claude Code} & Independent & 47 & 1.81$\times$ & 42 & 10 & 64 & 2.66$\times$ & 56 & 37 & 20 & 4.74$\times$ & 19 & 16 \\
 & Composed & 42 & 1.79$\times$ & 40 & 10 & 54 & 3.52$\times$ & 54 & 36 & 14 & 6.28$\times$ & 14 & 12 \\
 & Sequential & 42 & 1.76$\times$ & 40 & 10 & 57 & 3.21$\times$ & 57 & 37 & 20 & 4.50$\times$ & 20 & 15 \\
\midrule
\multirow{3}{*}{KDA} & Independent & 47 & 1.64$\times$ & 42 & 8 & 63 & 2.56$\times$ & 61 & 37 & 20 & 3.05$\times$ & 18 & 14 \\
 & Composed & 43 & 1.69$\times$ & 40 & 9 & 51 & 2.62$\times$ & 51 & 31 & 11 & 3.45$\times$ & 11 & 9 \\
 & Sequential & 43 & 1.68$\times$ & 40 & 8 & 57 & 2.30$\times$ & 57 & 30 & 20 & 2.86$\times$ & 20 & 14 \\
\midrule
\multirow{3}{*}{AKO} & Independent & 45 & 1.84$\times$ & 44 & 11 & 64 & 3.35$\times$ & 64 & 41 & 18 & 6.48$\times$ & 18 & 17 \\
 & Composed & 45 & 1.87$\times$ & 44 & 11 & 60 & 3.54$\times$ & 60 & 43 & 14 & 6.59$\times$ & 14 & 13 \\
 & Sequential & 45 & 1.86$\times$ & 44 & 11 & 59 & 3.51$\times$ & 59 & 40 & 20 & 6.55$\times$ & 20 & 17 \\
\midrule
\multirow{3}{*}{Codex} & Independent & 47 & 1.62$\times$ & 40 & 9 & 64 & 2.87$\times$ & 60 & 39 & 19 & 2.92$\times$ & 18 & 12 \\
 & Composed & 41 & 1.76$\times$ & 40 & 9 & 47 & 3.52$\times$ & 45 & 33 & 5 & 5.79$\times$ & 5 & 4 \\
 & Sequential & 41 & 1.75$\times$ & 38 & 9 & 56 & 2.90$\times$ & 56 & 34 & 18 & 3.95$\times$ & 18 & 11 \\
\bottomrule
\end{tabular}
}
\end{table}

\subsection{Agent Evaluation: End-to-End Leaderboard}
\label{sec:eval-e2e}

We deploy each agent's independent and sequential winner sets into the 11 default models (each receiving the kernels it instantiates) and run their real workloads on the production path, next to a baseline run and a noise run (Section~\ref{sec:metrics}).
If a set crashes the model or grossly corrupts its outputs, the offending kernel (named in the traceback, else found by bisection) is dropped and the model re-run; the first subset that runs correctly is the agent's \emph{deployable subset}, timed on the full workloads.
Table~\ref{tab:e2e-macroeval} scores both the full winner sets as-is (\emph{all winners}) and the deployable subsets (per model: Appendix~\ref{app:e2e}).

\begin{table}[t]
  \centering
  \caption{\textbf{End-to-end \textsc{MacroEval} leaderboard on the default set (B200).} \emph{Rank}: position of each agent--mode entry by Score (10 entries).
  \emph{Valid}: models with $C_i \geq \tau$; macro metrics as in Section~\ref{sec:metrics}; \emph{Thr.}/\emph{Lat.}: $\lambda{=}1$/$\lambda{=}0$.
  \emph{Deployable subset}: after dropping the kernels that crash or break the model; \emph{All winners}: every kernel-level winner as-is.}
  \label{tab:e2e-macroeval}
  \small
  \setlength{\tabcolsep}{3pt}
  \resizebox{\linewidth}{!}{% Generated by figures/make_macroeval.py (11/11 scenarios with a baseline, universe=all: 11 items, 8 families). Do not edit by hand.
% Rank column and bold best Score added by hand (ranks match tables/e2e_sensitivity.tex at 1x thresholds).
\begin{tabular}{@{}ll rrr rrr rr rr@{}}
\toprule
 & & \multicolumn{8}{c}{\textbf{Deployable subset}} & \multicolumn{2}{c}{\textbf{All winners}} \\
\cmidrule(lr){3-10}\cmidrule(lr){11-12}
 & & & & & \multicolumn{3}{c}{$S_{\mathrm{macro}}$} & & & & \\
\cmidrule(lr){6-8}
Agent & Mode & Valid & Cov$_{\mathrm{macro}}$ & $C_{\mathrm{macro}}$ & $\lambda{=}0.5$ & Thr. & Lat. & Score & Rank & Valid & Cov$_{\mathrm{macro}}$ \\
\midrule
\multirow{2}{*}{Dr.~Kernel} & Independent & 9/11 & 0.75 & 0.85 & 0.98$\times$ & 1.01$\times$ & 0.95$\times$ & 0.63 & 4 & 2/11 & 0.25 \\
 & Sequential & 9/11 & 0.75 & 0.75 & 0.99$\times$ & 1.02$\times$ & 0.96$\times$ & 0.55 & 5 & 2/11 & 0.25 \\
\midrule
\multirow{2}{*}{Claude Code} & Independent & 6/11 & 0.52 & 0.61 & 1.09$\times$ & 1.08$\times$ & 1.11$\times$ & 0.35 & 8 & 2/11 & 0.25 \\
 & Sequential & 7/11 & 0.56 & 0.55 & 1.11$\times$ & 1.09$\times$ & 1.14$\times$ & 0.34 & 9 & 1/11 & 0.12 \\
\midrule
\multirow{2}{*}{KDA} & Independent & 11/11 & 1.00 & 0.98 & 1.13$\times$ & 1.08$\times$ & 1.19$\times$ & \textbf{1.11} & \textbf{1} & 4/11 & 0.42 \\
 & Sequential & 11/11 & 1.00 & 0.98 & 1.04$\times$ & 1.03$\times$ & 1.05$\times$ & 1.03 & 2 & 6/11 & 0.58 \\
\midrule
\multirow{2}{*}{AKO} & Independent & 7/11 & 0.56 & 0.65 & 1.15$\times$ & 1.13$\times$ & 1.16$\times$ & 0.42 & 7 & 1/11 & 0.12 \\
 & Sequential & 9/11 & 0.83 & 0.82 & 1.25$\times$ & 1.08$\times$ & 1.44$\times$ & 0.85 & 3 & 2/11 & 0.25 \\
\midrule
\multirow{2}{*}{Codex} & Independent & 5/11 & 0.33 & 0.33 & 0.96$\times$ & 0.98$\times$ & 0.94$\times$ & 0.11 & 10 & 1/11 & 0.12 \\
 & Sequential & 8/11 & 0.71 & 0.71 & 1.10$\times$ & 1.11$\times$ & 1.08$\times$ & 0.55 & 6 & 1/11 & 0.12 \\
\bottomrule
\end{tabular}
}
\end{table}


\paragraph{Kernel-level winners rarely survive deployment as-is.}
Swapped in unchanged, only 22 of the 110 winner-set deployments run correctly (13--58\% of model families per agent).
Compilation-stack conflicts (e.g.\ embedding kernels calling raw C++ extensions that \texttt{torch.compile} cannot trace, rotary kernels that specialize dynamic batch shapes) and crashes on engine shapes the kernel bench never produced account for 86\% of the drops with an identified cause (Appendix Table~\ref{tab:e2e-drops}), all invisible to an isolated kernel check.
Other sets pass every kernel check yet corrupt the model (texture-only FLUX images, NaNs on Oasis), which only the end-to-end check flags.

\paragraph{Deployable end-to-end gains are small.}
Even after dropping the failing kernels, end-to-end speedups are modest ($S_{\mathrm{macro}}$ of $0.96$--$1.25\times$) despite kernel-level speedups of up to $6.6\times$: gains concentrate where the rewritten kernels dominate runtime (up to $2.2\times$ on YOLOv10, $1.7\times$ on OpenFold3, $1.5\times$ on GLA), while LLM serving stays within 6\% of the production kernels.
Sequential sets raise AKO and Codex from $1.15\times$ and $0.96\times$ to $1.25\times$ and $1.10\times$.

\paragraph{Kernel-level scores mis-rank agents.}
In isolation, Claude Code matches or beats KDA on correctness and speedup at every level (L3: 20 vs.\ 20 correct, $4.74\times$ vs.\ $3.05\times$), yet KDA's \textsc{MacroEval} score is $3\times$ higher ($1.11$ vs.\ $0.35$), the best of all agents even with halved or doubled thresholds (Appendix~\ref{app:e2e}).
The gap is deployability, not speed ($1.13\times$ vs.\ $1.09\times$): all 11 of KDA's models are valid versus 6 of Claude Code's. KDA wraps the same Claude model in a planning and review harness, which a kernel-level leaderboard cannot credit.

